% ---------------------------------------------------------------------
% Argomento: resistenze serie/parallelo, legge di Ohm, potenza
% Livello: base-intermedio   |   Tempo stimato: 12 min
% ---------------------------------------------------------------------
\begin{esercizio}
\dato{V}{12}{15}{18}
\dato{R1}{150}{180}{220}
\dato{R2}{220}{270}{330}
\dato{R3}{330}{560}{680}
\calcola[R_{23}]{R23}{\x{R2}*\x{R3}/(\x{R2}+\x{R3})}
\calcola[R_{eq}]{Req}{\x{R1}+\x{R23}}
\calcola[I]{I}{\x{V}/\x{Req}}
\calcola[P]{P}{\x{V}*\x{I}}

Dato il circuito in figura:

\begin{center}
\begin{minipage}{0.48\linewidth}\centering
  \begin{circuitikz}[scale=0.85, transform shape]
    \draw (0,0) to[vsource, l=$V$, n=G] (0,3)
          to[R, l=$R_1$, -*] (3,3) -- (5,3)
          (3,3) to[R, l_=$R_2$, -*] (3,0)
          (5,3) to[R, l=$R_3$] (5,0) -- (0,0);
    \morsetti{G}
  \end{circuitikz}
\end{minipage}\hfill
\begin{minipage}{0.4\linewidth}
  \renewcommand{\arraystretch}{1.3}
  \begin{tabular}{r@{${}={}$}l}
    $V$   & \mis{V}{\volt}\\
    $R_1$ & \mis{R1}{\ohm}\\
    $R_2$ & \mis{R2}{\ohm}\\
    $R_3$ & \mis{R3}{\ohm}
  \end{tabular}
\end{minipage}
\end{center}

\begin{parts}
  \part[2] Calcola la resistenza equivalente $R_{eq}$ vista dal generatore.
  \part[2] Calcola la corrente $I$ erogata dal generatore.
  \part[2] Calcola la potenza $P$ erogata dal generatore.
\end{parts}

\begin{solution}
  $R_{23} = \dfrac{R_2 R_3}{R_2+R_3} = \dfrac{\val{R2}\cdot\val{R3}}{\val{R2}+\val{R3}}
    \approx \mis{R23}{\ohm}$,\qquad
  $R_{eq} = R_1 + R_{23} \approx \mis{Req}{\ohm}$\\[4pt]
  $I = V/R_{eq} \approx \mis{I}{\ampere}$\qquad
  $P = V\cdot I \approx \mis{P}{\watt}$
  \criterio{a) 1 pt riconoscere il parallelo $R_2\parallel R_3$, 1 pt calcolo;
    b) e c) 1 pt formula, 1 pt risultato con unità.}
\end{solution}
\end{esercizio}
